%% ============================================================
%%  Abstract
%% ============================================================
\chapter*{Abstract}
\addcontentsline{toc}{chapter}{Abstract}

Modern networked environments face an increasingly sophisticated threat landscape
in which adversaries pursue structured campaigns across multiple exploitation
phases.
Traditional intrusion-detection and honeypot management systems rely on static
rule bases or fixed-threshold anomaly detectors that cannot adapt as attacker
strategies evolve in real time.
This thesis presents \textbf{HoneyIQ}, a simulation framework and decision
policy for adaptive honeypot management grounded in the Lockheed~Martin
Cyber Kill Chain~\citep{hutchins2011intelligence}.
This revised edition extends the original framework with session-coherent
synthetic traffic, a severity-weighted escalation signal, and two auditable
dynamic-response mechanisms, and --- motivated directly by scrutiny of the
original evaluation protocol --- replaces bare point-estimate headline
metrics with a mixed-traffic protocol reported alongside confidence
intervals and raw sample counts.

The attacker subsystem models adversarial behaviour as a pair of coupled
Discrete-Time Markov Chains governing attack-type and kill-chain-stage
transitions.
Network traffic features are generated by sampling from parametric
distributions calibrated against the UNSW-NB15 intrusion detection
benchmark~\citep{moustafa2015unsw}.
Four attacker intent profiles --- \emph{Stealthy}, \emph{Aggressive},
\emph{Targeted}, and \emph{Opportunistic} --- are encoded as intent-specific
transition modifiers, producing qualitatively distinct attack campaigns from
a shared base model.
In this revision, each simulated session additionally draws a single,
persistent intensity factor and --- for benign traffic --- one of three
weighted personas (\emph{casual user}, \emph{crawler}, \emph{monitoring
probe}), replacing per-step independent resampling with session-coherent
traffic and heterogeneous benign behaviour; the supporting classifier was
retrained on this richer generator.

The principal contribution of this thesis is the
\textbf{Stage-Escalation Decision Matrix (SEDM)}, a transparent, deterministic
policy that maps the current kill chain stage and a Markov-chain-derived
escalation risk score to one of five honeypot responses
(ALLOW, LOG, TROLL, BLOCK, ALERT).
A composite threat-level signal integrating attack severity, kill chain
progress, escalation rate, and cumulative attack pressure provides a
single normalised index for situational awareness.
A Random Forest classifier~\citep{breiman2001random} provides attack-type
identification from raw network features as a supporting component.
This revision extends the SEDM in three additive, auditable directions: a
severity-weighted exponential moving average as an alternative to the
original hard-window escalation signal; a half-life-decayed, cross-session
reputation score with a new override rule (R4) that persists suspicion
across a source's individual session boundaries, fail2ban-style; and a
bounded deadband controller that nudges the rate-override threshold to
curb alert fatigue, explicitly scoped as a frequency-management safety
valve rather than a correctness mechanism, since no ground-truth feedback
signal exists in this system to justify a stronger claim.

Experimental evaluation across 30~episodes per attacker intent (200 steps
each, continuous-attack protocol) demonstrates that the SEDM achieves
consistently high threat coverage: detection rates of \textbf{99.93\%},
\textbf{100.0\%}, \textbf{99.97\%}, and \textbf{99.97\%} for the Stealthy,
Aggressive, Targeted, and Opportunistic profiles respectively, with a
measured false positive rate of \textbf{0.00\%} across all four intents
under classifier-driven decisions.
This near-zero false-positive rate is not a general immunity claim: it is a
direct consequence of the classifier's near-perfect separation of the
synthetic attack-type distributions under a protocol containing almost no
benign traffic to misclassify, and both a feature-noise robustness sweep
(false positives rising to 10.00\% at 50\% injected feature noise) and a
revised mixed-traffic protocol --- 50 episodes per intent, 500 steps each,
30\% benign-traffic injection, with detection and false-positive rates
reported as Wilson 95\% confidence intervals over their full raw sample
counts rather than as bare percentages --- confirm that the headline
zero-FPR figure is a property of the original protocol's traffic
composition rather than of the policy itself.
The traffic-realism extensions alone measurably, if modestly, widen this
gap (classifier NORMAL-class recall falls from 1.0000 to 0.9900; end-to-end
specificity under mixed traffic falls from 0.999 to 0.992--0.995), evidence
generated without any artificial noise injection.
The SEDM's transparency, zero training overhead, and consistent cross-intent
performance are compared with the Deep Q-Network (DQN) baseline, whose
training on Opportunistic traffic reproduces a detection-rate improvement
but converges to a degenerate near-always-escalate policy (false-positive
rate pinned near 1.0) rather than genuine discrimination; this outcome,
together with two previously undetected label-alignment defects in the
evaluation harness itself, is presented as direct evidence against adopting
a reward-learned mechanism for dynamic response in place of the auditable
reputation and threshold-adaptation mechanisms introduced here.

\vspace{1em}
\noindent\textbf{Keywords:} cyber kill chain, honeypot management, Stage-Escalation
Decision Matrix, deep reinforcement learning, Deep Q-Network, Markov chain,
intrusion detection, attacker intent, UNSW-NB15, adaptive cybersecurity,
cross-session reputation, confidence intervals.
\clearpage
